\chapter{Introducción}
\label{sec.ej3:introduccion}

La modulación de amplitud constituye uno de los pilares históricos y conceptuales en los sistemas de comunicaciones
analógicas. El proceso consiste en variar la amplitud instantánea de una portadora de alta frecuencia de forma
proporcional a la señal de información en banda base. Dependiendo de si se conserva o no la componente espectral de la
portadora en la transmisión, se originan dos variantes principales: la modulación en doble banda lateral con portadora
completa (AM estándar o DSB-FC) y la modulación en doble banda lateral con portadora suprimida (DSB-SC).

En este trabajo práctico se analiza el comportamiento experimental de ambos esquemas de modulación y demodulación
utilizando el circuito integrado modulador balanceado MC1496.

\section{Objetivos del Ensayo}
Los objetivos fijados para el desarrollo de esta experiencia de laboratorio son:
\begin{itemize}
    \item Implementar experimentalmente los esquemas de modulación y demodulación de señales DSB-SC y AM estándar
          empleando el Kit de Comunicaciones N° 5 basado en el integrado MC1496.
    \item Observar y caracterizar en el dominio del tiempo las formas de onda correspondientes a la señal modulante,
          la portadora de radiofrecuencia y la señal resultante a la salida del modulador.
    \item Contrastar las mediciones temporales con el análisis espectral en frecuencia mediante la transformada rápida
          de Fourier (FFT) integrada en el osciloscopio digital.
    \item Evaluar el efecto de la variación de amplitud y frecuencia de la banda base sobre el índice de modulación,
          el ancho de banda ocupado y los niveles de potencia de las bandas laterales.
    \item Caracterizar el fenómeno de sobremodulación en AM y constatar la generación de productos de distorsión
          espectral.
    \item Analizar el proceso de demodulación sincrónica (detector de producto) y el filtrado paso bajo posterior.
    \item Medir y cuantificar la capacidad de rechazo y balance de portadora en el modulador balanceado a través
          del ajuste de preset, expresando la atenuación alcanzada en decibeles.
\end{itemize}

\section{Equipamiento e Instrumental Utilizado}
Para la realización de los ensayos de laboratorio se dispuso del siguiente instrumental:
\begin{itemize}
    \item Kit de Comunicaciones N° 5 de la Cátedra (módulo de modulación y detección de producto con dos integrados MC1496).
    \item Osciloscopio digital Rigol de dos canales con capacidad de cálculo de FFT en tiempo real.
    \item Dos generadores de señales de funciones senoidales calibrados (para banda base y portadora).
    \item Fuente de alimentación regulada de corriente continua con salidas de $+12\V$ y $-8\V$.
    \item Puntas de prueba de osciloscopio atenuadas $10\times$ con terminal de referencia a tierra.
\end{itemize}

